\hypertarget{python-1.0-to-1.6-inception-the-ll1-executable-pipeline}{%
\chapter{Python 1.0 to 1.6: Inception \& the LL(1) Executable
Pipeline}\label{python-1.0-to-1.6-inception-the-ll1-executable-pipeline}}

\hypertarget{chronological-inception-core-design-philosophies}{%
\subsection{1.1 Chronological Inception \& Core Design
Philosophies}\label{chronological-inception-core-design-philosophies}}

Python was conceived in December 1989 by Guido van Rossum at CWI
(Centrum Wiskunde \& Informatica) in the Netherlands. It was created as
a successor to the \textbf{ABC language}, a programming language
designed for beginners that was elegant but suffered from fatal
architectural flaws: 1. \textbf{Monolithic Runtime}: ABC was difficult
to extend. It did not support dynamic library loading or simple
integration with C. 2. \textbf{Strict Filesystem Database}: ABC stored
variable states globally in a monolithic filesystem database, which
severely limited performance and concurrent operations. 3.
\textbf{Closed Ecosystem}: ABC was a closed system, making it impossible
to interface directly with low-level OS structures or hardware calls.

Van Rossum designed Python as an open-source, extensible scripting
language that addressed these issues: * \textbf{Whitespace
Significance}: Unlike languages that use curly braces
(\passthrough{\lstinline!\{\}!}), CPython's tokenizer tracks
indentations (\passthrough{\lstinline!INDENT!}) and dedentations
(\passthrough{\lstinline!DEDENT!}) as formal syntax control elements.
This ensures visual block structure matches execution scopes. *
\textbf{Unified Object Model}: Every data structure is a heap-allocated
object (\passthrough{\lstinline!PyObject!}), allowing unified interfaces
and straightforward extensibility via C extensions. * \textbf{Names and
Bindings}: Python variables are not declared memory slots containing
values; they are references (pointers) to heap-allocated objects. Thus,
\passthrough{\lstinline!x = 5!} binds the label
\passthrough{\lstinline!x!} in the namespace dictionary to the
\passthrough{\lstinline!PyObject!} representing the integer 5.

\begin{center}\rule{0.5\linewidth}{0.5pt}\end{center}

\hypertarget{cpython-parser-evolution-ll1-left-recursion-math-vs.-peg}{%
\subsection{1.2 CPython Parser Evolution: LL(1) Left-Recursion Math
vs.~PEG}\label{cpython-parser-evolution-ll1-left-recursion-math-vs.-peg}}

Until Python 3.9, CPython used a custom \textbf{LL(1) parser}
(Left-to-right, Leftmost derivation with 1 token lookahead) to construct
its concrete parse trees.

\hypertarget{left-recursion-in-ll1}{%
\subsubsection{1. Left-Recursion in LL(1)}\label{left-recursion-in-ll1}}

An LL(1) parser is a top-down parser. It cannot parse grammar rules that
exhibit \textbf{left-recursion}. A grammar rule is left-recursive if the
leftmost symbol on the right-hand side of a production is the
non-terminal itself: \[A \to A \alpha \mid \beta\]

When a top-down parser attempts to expand \(A\) using this rule, it
loops infinitely without consuming any input:
\[A \to A \alpha \to A \alpha \alpha \to A \alpha \alpha \alpha \to \dots\]

\hypertarget{grammar-rewriting-workarounds}{%
\subsubsection{2. Grammar Rewriting
Workarounds}\label{grammar-rewriting-workarounds}}

To prevent infinite recursion, developers had to rewrite left-recursive
rules into right-recursive forms: \[A \to \beta A'\]
\[A' \to \alpha A' \mid \epsilon\]

While mathematically equivalent, this workaround made the compiler's
grammar definition complex, hard to maintain, and less readable.

\hypertarget{parsing-expression-grammar-peg-parser-pep-617}{%
\subsubsection{3. Parsing Expression Grammar (PEG) Parser (PEP
617)}\label{parsing-expression-grammar-peg-parser-pep-617}}

Python 3.9 replaced the LL(1) parser with a \textbf{PEG parser}. *
\textbf{Packrat Parsing (Memoization)}: A PEG parser handles
left-recursion by tracking and caching parsed rules at each input index.
If the parser backtracks or encounters left-recursion, it references the
cache to break infinite loops. * \textbf{Infinite Lookahead}: Unlike
LL(1)'s single-token limit, the PEG parser can look ahead arbitrarily
far, allowing more expressive syntax. * \textbf{Direct AST Generation}:
The PEG parser generates the Abstract Syntax Tree (AST) directly from
the tokenizer stream, eliminating the overhead of building an
intermediate Concrete Parse Tree.

\begin{lstlisting}
LL(1) Parser Pipeline (Classic):
[Source] -> [Tokenizer] -> [Concrete Parse Tree] -> [AST Generator] -> [AST]

PEG Parser Pipeline (Modern 3.9+):
[Source] -> [Tokenizer] -> [PEG Parser (Memoized)] -> [AST]
\end{lstlisting}

\begin{center}\rule{0.5\linewidth}{0.5pt}\end{center}

\hypertarget{compilation-ast-and-the-symbol-table-pass}{%
\subsection{1.3 Compilation, AST, and the Symbol Table
Pass}\label{compilation-ast-and-the-symbol-table-pass}}

After generating the AST, CPython compiles the code. This begins with a
static analysis pass to build the symbol table
(\passthrough{\lstinline!Python/symtable.c!}).

\hypertarget{variable-scope-classification}{%
\subsubsection{1. Variable Scope
Classification}\label{variable-scope-classification}}

CPython classifies variable names into four distinct scopes: 1.
\textbf{Local}: Bound within the current function scope. 2.
\textbf{Global Explicit}: Declared via
\passthrough{\lstinline!global x!}. 3. \textbf{Global Implicit}:
Referenced but not bound within the current function (resolved at
runtime in module globals). 4. \textbf{Enclosing (Closure/Free)}:
Variables defined in an outer function scope and accessed by a nested
inner function.

\hypertarget{closure-cells-pycellobject}{%
\subsubsection{\texorpdfstring{2. Closure Cells
(\texttt{PyCellObject})}{2. Closure Cells (PyCellObject)}}\label{closure-cells-pycellobject}}

When a local variable is accessed by a nested function, it cannot be
stored in the standard fast-locals array of the outer function because
its lifetime must extend beyond the outer function's execution. CPython
handles this by wrapping the variable in a
\textbf{\passthrough{\lstinline!PyCellObject!}}.

\begin{lstlisting}[language=C]
typedef struct {
    PyObject_HEAD
    PyObject *ob_ref;       /* Pointer to the enclosed PyObject */
} PyCellObject;
\end{lstlisting}

The cell object resides on the heap, allowing both the outer and inner
functions to share access to the same object reference, even after the
outer function's stack frame is popped.

\hypertarget{the-dynamic-namespace-trap}{%
\subsubsection{3. The Dynamic Namespace
Trap}\label{the-dynamic-namespace-trap}}

Python allows dynamic namespace modifications using features like
\passthrough{\lstinline!exec()!}, \passthrough{\lstinline!eval()!}, or
\passthrough{\lstinline!from module import *!}. When the compiler
detects these dynamic features inside a function, it cannot determine
variable scopes at compile-time. As a result, it disables fast local
lookups: * \textbf{Standard Local Lookup}: Uses the
\passthrough{\lstinline!LOAD\_FAST!} bytecode, which retrieves
references from an array index in \(\mathcal{O}(1)\) time. *
\textbf{Dynamic Fallback}: Reverts to the slow
\passthrough{\lstinline!LOAD\_NAME!} bytecode, which searches the local,
enclosing, global, and builtin dictionaries sequentially at runtime.

\hypertarget{symbol-table-inspection}{%
\subsubsection{4. Symbol Table
Inspection}\label{symbol-table-inspection}}

We can inspect this static compilation analysis using the built-in
\passthrough{\lstinline!symtable!} module:

\begin{lstlisting}[language=Python]
import symtable

code_text = """
def outer_func(x):
    z = 10
    def inner_func():
        nonlocal z
        return x + z
    return inner_func
"""

table = symtable.symtable(code_text, filename="<string>", compile_type="exec")
outer_scope = table.lookup("outer_func").get_namespace()

print("Outer variables:", outer_scope.get_identifiers())
print("Is 'z' a cell variable in outer_func?", outer_scope.lookup("z").is_cell())
print("Is 'x' a cell variable in outer_func?", outer_scope.lookup("x").is_cell())

inner_scope = outer_scope.lookup("inner_func").get_namespace()
print("Inner variables:", inner_scope.get_identifiers())
print("Is 'z' a free (closure) variable in inner_func?", inner_scope.lookup("z").is_free())
\end{lstlisting}

\begin{center}\rule{0.5\linewidth}{0.5pt}\end{center}

\hypertarget{pycodeobject-bytecode-anatomy}{%
\subsection{1.4 PyCodeObject \& Bytecode
Anatomy}\label{pycodeobject-bytecode-anatomy}}

The compiler processes the AST and symbol table to generate a
\passthrough{\lstinline!PyCodeObject!} struct, which serves as the
immutable blueprint for execution.

\hypertarget{c-level-structure-of-pycodeobject}{%
\subsubsection{\texorpdfstring{1. C-Level Structure of
\texttt{PyCodeObject}}{1. C-Level Structure of PyCodeObject}}\label{c-level-structure-of-pycodeobject}}

Defined in \passthrough{\lstinline!Include/cpython/code.h!}, this
structure holds the compiled bytecode, constants pool, and variable name
catalogs:

\begin{lstlisting}[language=C]
struct PyCodeObject {
    PyObject_HEAD
    int co_argcount;            /* Number of positional arguments */
    int co_posonlyargcount;     /* Positional-only arguments (Python 3.8+) */
    int co_kwonlyargcount;      /* Keyword-only arguments */
    int co_nlocals;             /* Number of local variables */
    int co_stacksize;           /* Maximum value stack depth needed by VM */
    int co_flags;               /* Compiler configuration flags (e.g. CO_GENERATOR) */
    
    PyObject *co_code;          /* Instruction sequence (bytes or 16-bit code units) */
    PyObject *co_consts;        /* Immutable literal constants (tuple of PyObjects) */
    PyObject *co_names;         /* Non-local/Global names referenced (tuple of strings) */
    PyObject *co_varnames;      /* Local parameter and variable names (tuple) */
    PyObject *co_freevars;      /* Free variables captured in closures (tuple) */
    PyObject *co_cellvars;      /* Local variables enclosed by nested scopes (tuple) */
};
typedef struct PyCodeObject PyCodeObject;
\end{lstlisting}

\hypertarget{instruction-decoding-and-extended_arg}{%
\subsubsection{\texorpdfstring{2. Instruction Decoding and
\texttt{EXTENDED\_ARG}}{2. Instruction Decoding and EXTENDED\_ARG}}\label{instruction-decoding-and-extended_arg}}

\begin{itemize}
\tightlist
\item
  \textbf{Bytecode Formats}:

  \begin{itemize}
  \tightlist
  \item
    \emph{Pre-3.11}: Instructions used a 2-byte format (1 byte for the
    opcode, 1 byte for the argument).
  \item
    \emph{Python 3.11+}: Instructions use 16-bit code units (2 bytes),
    combining opcode, arguments, and inline caches.
  \end{itemize}
\item
  \textbf{The \passthrough{\lstinline!EXTENDED\_ARG!} Opcode}: The
  argument field of a standard instruction is limited to 8 bits
  (representing indices \(0\) to \(255\)). If an index exceeds 255
  (e.g., loading the 300th constant in
  \passthrough{\lstinline!co\_consts!}), the compiler prefixes the
  instruction with an \passthrough{\lstinline!EXTENDED\_ARG!} opcode.

  \begin{itemize}
  \tightlist
  \item
    The \passthrough{\lstinline!EXTENDED\_ARG!} instruction loads the
    high-order bits of the argument value into an internal register.
  \item
    The subsequent instruction shifts this value and adds its own
    argument payload, allowing the VM to resolve indices up to 32 bits
    wide.
  \end{itemize}
\end{itemize}

Let's write a script to inspect these low-level compiled fields inside a
live Python runtime:

\begin{lstlisting}[language=Python]
def example_fn(a, b):
    local_val = 42
    return a + b + local_val

code_obj = example_fn.__code__

print("Local Variable Names (co_varnames):", code_obj.co_varnames)
print("Constant Pool (co_consts):", code_obj.co_consts)
print("Instruction Byte Sequence (co_code):", list(code_obj.co_code))
print("Scoping/Compiler Flags (co_flags):", bin(code_obj.co_flags))
\end{lstlisting}

\begin{center}\rule{0.5\linewidth}{0.5pt}\end{center}

\hypertarget{the-virtual-machine-execution-loop-frame-objects}{%
\subsection{1.5 The Virtual Machine Execution Loop \& Frame
Objects}\label{the-virtual-machine-execution-loop-frame-objects}}

When a function is called, the CPython virtual machine allocates a new
stack frame representation, the \passthrough{\lstinline!PyFrameObject!},
on the heap.

\hypertarget{c-level-structure-of-pyframeobject}{%
\subsubsection{\texorpdfstring{1. C-Level Structure of
\texttt{PyFrameObject}}{1. C-Level Structure of PyFrameObject}}\label{c-level-structure-of-pyframeobject}}

Defined in \passthrough{\lstinline!Include/internal/pycore\_frame.h!},
the frame object tracks runtime execution state:

\begin{lstlisting}[language=C]
struct _frame {
    PyObject_HEAD
    struct _frame *f_back;      /* Link to previous frame (caller's frame) */
    PyCodeObject *f_code;       /* Code object executed in this frame */
    PyObject *f_builtins;       /* Builtins symbol lookup dictionary */
    PyObject *f_globals;        /* Global namespace dictionary */
    PyObject *f_locals;         /* Local namespace dictionary */
    PyObject **f_valuestack;    /* Points to the base of the evaluation stack */
    int f_lasti;                /* Last instruction index executed */
    /* ... additional traceback and debugging fields ... */
};
typedef struct _frame PyFrameObject;
\end{lstlisting}

\hypertarget{cpython-vm-evaluation-loop-_pyeval_evalframedefault}{%
\subsubsection{\texorpdfstring{2. CPython VM Evaluation Loop
(\texttt{\_PyEval\_EvalFrameDefault})}{2. CPython VM Evaluation Loop (\_PyEval\_EvalFrameDefault)}}\label{cpython-vm-evaluation-loop-_pyeval_evalframedefault}}

The core execution engine is defined in
\passthrough{\lstinline!Python/ceval.c!} inside the function
\passthrough{\lstinline!\_PyEval\_EvalFrameDefault()!}. This function
contains a monolithic evaluation loop:

\begin{lstlisting}[language=C]
PyObject* _PyEval_EvalFrameDefault(PyThreadState *tstate, PyFrameObject *f, int throwflag) {
    // Local optimization pointers
    PyObject **stack_pointer = f->f_valuestack;
    const _Py_CODEUNIT *next_instr = (_Py_CODEUNIT *)PyBytes_AS_STRING(f->f_code->co_code);
    
    for (;;) {
        _Py_CODEUNIT word = *next_instr++;
        int opcode = _Py_OPCODE(word);
        int oparg = _Py_OPARG(word);
        
        switch (opcode) {
            case LOAD_FAST: {
                PyObject *value = GETLOCAL(oparg);
                Py_INCREF(value);
                PUSH(value);
                DISPATCH();
            }
            case BINARY_OP: {
                PyObject *right = POP();
                PyObject *left = POP();
                PyObject *res = PyNumber_Add(left, right);
                Py_DECREF(left);
                Py_DECREF(right);
                PUSH(res);
                DISPATCH();
            }
            case RETURN_VALUE: {
                PyObject *retval = POP();
                return retval;
            }
            /* ... additional opcodes ... */
        }
    }
}
\end{lstlisting}

\hypertarget{value-stack-execution-walkthrough}{%
\subsubsection{3. Value Stack Execution
Walkthrough}\label{value-stack-execution-walkthrough}}

CPython is a stack-based virtual machine. Operations pop arguments from
the evaluation stack and push results back. Here is a trace of
evaluating the expression \passthrough{\lstinline!a + b!}:

\begin{lstlisting}
Stack State Trace for BINARY_OP (Add):

1. LOAD_FAST 0 (a)     2. LOAD_FAST 1 (b)     3. BINARY_OP           4. RETURN_VALUE
   +-----------+          +-----------+          +-----------+          +-----------+
   |   val_a   |          |   val_b   |          |  val_a+b  |          |  (Empty)  |
   +-----------+          +-----------+          +-----------+          +-----------+
   |  (Empty)  | -->      |   val_a   | -->      |  (Empty)  | -->      |  (Empty)  |
   +-----------+          +-----------+          +-----------+          +-----------+
\end{lstlisting}

The program counter (\passthrough{\lstinline!next\_instr!}) increments,
fetching code units, popping parameters off the stack, and executing the
corresponding C functions.

\begin{center}\rule{0.5\linewidth}{0.5pt}\end{center}
